\subsection{Perkalian Skalar Vektor dalam \texorpdfstring{$\mathbb{R}^n$}{Rn}}

Perkalian skalar vektor dalam $\mathbb{R}^n$ didefinisikan sebagai \index{vektor!perkalian skalar}
berikut.

\begin{definition}{Perkalian Skalar Vektor dalam $\mathbb{R}^n$}\label{def:vectorscalarmultiplication}
Jika $\vect{u}\in \mathbb{R}^{n}$ dan $k\in \mathbb{R}$ adalah suatu
skalar,
\index{skalar} maka $k\vect{u}\in \mathbb{R}^{n}$ didefinisikan oleh
\begin{equation*}
k\vect{u}=k\leftB \begin{array}{c}
u_{1} \\
\vdots \\
u_{n}
\end{array}
\rightB = \leftB \begin{array}{c}
ku_{1} \\
\vdots \\
ku_{n}
\end{array}
\rightB
\end{equation*}
\end{definition}

Seperti halnya penjumlahan, perkalian skalar vektor memenuhi beberapa sifat penting. Sifat-sifat ini
diuraikan dalam teorema berikut.

\begin{theorem}{Sifat-Sifat Perkalian Skalar}\label{thm:vectorscalarmult}
Sifat-sifat berikut berlaku untuk vektor $\vect{u},\vect{v}\in \mathbb{R}^{n}$ dan skalar $k,p $.
\begin{itemize}
\item Hukum Distributif terhadap Penjumlahan Vektor
\begin{equation*}
k \left( \vect{u}+\vect{v}\right) = k\vect{u}+ k\vect{v}
\end{equation*}
\item Hukum Distributif terhadap Penjumlahan Skalar
\begin{equation*}
\left( k + p  \right)\vect{u} = k \vect{u}+p \vect{u}
\end{equation*}
\item Hukum Asosiatif untuk Perkalian Skalar
\begin{equation*}
k \left( p \vect{u}\right) = \left(k p \right)\vect{u}
\end{equation*}
\item Aturan Perkalian dengan $1$
\begin{equation*}
1\vect{u}=\vect{u}
\end{equation*}
\end{itemize}
\end{theorem}

\ifdefined\showproofs
\begin{proof}
Kita akan menunjukkan bukti untuk:
\begin{equation*}
k \left( \vect{u}+\vect{v}\right) = k \vect{u}+ k \vect{v}
\end{equation*}
Perhatikan bahwa:
\begin{equation*}
\begin{array}{ll}
k \left( \vect{u}+\vect{v}\right) & =k \leftB u_{1}+v_{1} \cdots u_{n}+v_{n}\rightB^T \\
& = \leftB k \left( u_{1}+v_{1}\right) \cdots k \left( u_{n}+v_{n}\right) \rightB^T \\
& = \leftB k u_{1}+ k  v_{1} \cdots k u_{n}+ k v_{n}\rightB^T \\
& = \leftB k u_{1} \cdots k u_{n} \rightB^T + \leftB k v_{1} \cdots k v_{n} \rightB^T \\
& = k \vect{u}+k \vect{v} \\
\end{array}
\end{equation*}
\end{proof}
\fi
Sekarang kita sajikan sebuah gagasan berguna yang mungkin pernah Anda jumpai sebelumnya, yang menggabungkan penjumlahan vektor dan perkalian skalar.

\begin{definition}{Kombinasi Linear}\label{def:linearcombination}
Sebuah vektor $\vect{v}$ disebut sebagai \textbf{kombinasi linear}
\index{kombinasi linear} dari vektor-vektor $\vect{u}_1,\cdots , \vect{u}_n $
jika terdapat skalar $a_{1},\cdots ,a_{n}$ sedemikian sehingga
\begin{equation*}
\vect{v} = a_1 \vect{u}_1 + \cdots + a_n \vect{u}_n
\end{equation*}
\end{definition}

Sebagai contoh,
\begin{equation*}
3
\leftB
\begin{array}{r}
-4 \\
1 \\
0
\end{array}
\rightB
+
2
\leftB
\begin{array}{r}
-3 \\
0\\
1
\end{array}
\rightB
 =
\leftB
\begin{array}{r}
-18 \\
3 \\
2
\end{array}
\rightB.
\end{equation*}
Dengan demikian, kita dapat mengatakan bahwa
\begin{equation*}
\vect{v}= \leftB
\begin{array}{r}
-18 \\
3 \\
2
\end{array}
\rightB
\end{equation*}
merupakan kombinasi linear dari vektor-vektor
\begin{equation*}
\vect{u}_1 = \leftB
\begin{array}{r}
-4 \\
1 \\
0
\end{array}
\rightB
\mbox{ dan }
\vect{u}_2 =
\leftB
\begin{array}{r}
-3 \\
0\\
1
\end{array}
\rightB
\end{equation*}
